\documentclass[10pt,a4paper]{amsart}
\usepackage[margin=19mm]{geometry}
\usepackage{amsmath,amssymb,mathtools}
\usepackage[scr=rsfs,cal=euler]{mathalfa}
\usepackage{xfrac}
\usepackage[unicode,hidelinks,bookmarks=false]{hyperref}
\newcommand{\ie}{i.e.\ }
\newcommand{\cf}{cf.\ }
\newcommand{\resp}{respectively\ }
\newcommand{\Subsection}{\subsection}
\providecommand{\nobreakdash}{\nobreak\hbox{-}}
\setlength{\emergencystretch}{2em}
\multlinegap0pt
\usepackage{bm}
\usepackage{bbm}
\usepackage{dsfont}
\usepackage{xspace}
\usepackage{cancel}
\usepackage{mathtools}
\usepackage{cleveref}
\usepackage[shortlabels]{enumitem}
\usepackage{amsthm}
\makeatletter
\@ifundefined{theorem}{\newtheorem{theorem}{Theorem}}{}
\makeatother
\newcommand{\ZZ}{\mathbb Z}
\title[Critical-exponent spectra and rank two inverse realization o]{Critical-exponent spectra and rank two inverse realization on biregular trees}
\author{Sanghoon Kwon}
\date{Source: arXiv:2607.21294v2 (CC BY 4.0)}
\begin{document}
\maketitle

\noindent\emph{Source and licence.} Critical-exponent spectra and rank two inverse realization on biregular trees. Version arXiv:2607.21294v2 (CC BY 4.0), distributed under the Creative Commons Attribution 4.0 International licence, as recorded in the arXiv licence metadata for this exact version and shown on the abstract page https://arxiv.org/abs/2607.21294v2. Full rights notice: licenses/arxiv-2607.21294-CC-BY-4.0.txt.\par\smallskip

\noindent\textit{Editorial scope.} Counted unit: the Theorem whose source label is \texttt{thm:rank-two-polynomials} in the article (source0.tex lines 2030--2074, source label \texttt{thm:rank-two-polynomials}), delivered together with the article's own complete original proof of it (source0.tex lines 2076--2139, one \texttt{proof} environment of 64 lines). No mathematics is written, repaired or re-derived: statement and proof are the article's own text line for line. Required context retained uncounted: thm:kernel-pressure (lines 1624--1655). Every definition, display and label of those lines is retained unchanged. Omitted from this package and not redistributed: 1--1623, 1656--2029, 2075--2075, 2140--2622. Nothing inside a retained excerpt is omitted or altered: every retained body is delivered exactly as the article's lines are written, so no omission receipt of any kind accompanies this package. Citations: the retained excerpts carry no citation command at all, so no bibliography entry of the article is transcribed and no reference is redistributed. Reference adaptation: none was needed and none was made; every reference inside the delivered unit resolves inside it, so no unresolved reference or citation is produced. Printed slips kept literal: no printed word, symbol, hypothesis or number of the counted statement or of its proof was altered. Layout and class: the article's own class is \texttt{amsart} with packages including geometry, amsmath, amssymb, amsthm, mathtools, enumitem, float, microtype, placeins, lineno, xcolor, tikz, hyperref, cleveref; the wrapper is the lane's pinned \texttt{amsart} editorial wrapper at 10pt on A4, which restates the theorem-like environments the retained text needs and adds the article's own macro definitions that the retained text needs (\texttt{\textbackslash ZZ}), with extra wrapper packages (bm, bbm, dsfont, xspace, cancel, mathtools, cleveref, enumitem). The delivered numbering is the unit's own. Assets: no retained span contains a figure environment, a graphics-inclusion command, a PDF-page inclusion, a table or a TikZ picture; no image file is copied into this package, the package declares no asset dependency and no image file is redistributed with it. Licence: version arXiv:2607.21294v2 of this article is distributed under the Creative Commons Attribution 4.0 International licence, as recorded in the arXiv licence metadata for this exact version and shown on the abstract page https://arxiv.org/abs/2607.21294v2. A full rights notice is delivered at licenses/arxiv-2607.21294-CC-BY-4.0.txt. Characters: every retained excerpt is pure ASCII, so the delivered engine, XeLaTeX with its default Latin Modern text font, reports no missing character.
\begin{theorem}[Kernel pressure equation]\label{thm:kernel-pressure}
Let \(C=K(\ell)\) be a finite connected non-cyclic core.  There is a unique
number \(\alpha_{K,\ell}>1\) such that
\begin{equation}\label{eq:kernel-pressure-defined}
  \rho\!\left(M_{K,\ell}(\alpha_{K,\ell})\right)=1,
\end{equation}
and this number is
\[
  \alpha_{K,\ell}=\rho(B_C).
\]
Equivalently, for every prescribed \(\alpha>1\),
\begin{equation}\label{eq:kernel-pressure}
  \rho(B_C)=\alpha
  \quad\Longleftrightarrow\quad
  \rho\!\left(M_{K,\ell}(\alpha)\right)=1.
\end{equation}
The function \(x\mapsto\rho(M_{K,\ell}(x))\) is continuous and strictly
decreasing from a value greater than one to zero.

The pressure-selected number \(\alpha_{K,\ell}\) is the distinguished
positive zero of the monic integral polynomial
\begin{equation}\label{eq:kernel-polynomial}
  F_{K,\ell}(x)
  =\det\!\left(D_\ell(x)-T_K\right)\in\ZZ[x].
\end{equation}
Here ``distinguished'' means that
\(\rho(M_{K,\ell}(\alpha_{K,\ell}))=1\); no uniqueness assertion is made
about positive zeros of the determinant that are not selected by the
spectral-radius equation.  In particular, the minimal polynomial of
\(\alpha_{K,\ell}\) divides
\(F_{K,\ell}(x)\).
\end{theorem}

\begin{theorem}[Rank-two polynomial classification]
\label{thm:rank-two-polynomials}
Let \(C\) be a finite connected core with \(b_1(C)=2\), and put
\(\alpha=\rho(B_C)\).  The kernel of \(C\) is of exactly one of the following
three types.
\begin{enumerate}[label=\textup{(\roman*)}]
  \item If \(C\) is a figure-eight with cycle lengths \(a,b\geq1\), then
  \begin{align}
    \frac{2}{\alpha^a+1}+\frac{2}{\alpha^b+1}&=1,
      \label{eq:figure-eight-pressure}\\
    (\alpha^a-1)(\alpha^b-1)&=4,                \label{eq:figure-eight-product}
  \end{align}
  and \(\alpha\) is the unique root \(x>1\) of
  \begin{equation}\label{eq:figure-eight-polynomial}
    R_{a,b}(x)=x^{a+b}-x^a-x^b-3.
  \end{equation}

  \item If \(C\) is a theta graph with path lengths \(a,b,c\geq1\), then
  \begin{equation}\label{eq:theta-pressure}
    \frac{1}{\alpha^a+1}
    +\frac{1}{\alpha^b+1}
    +\frac{1}{\alpha^c+1}=1,
  \end{equation}
  and \(\alpha\) is the unique root \(x>1\) of
  \begin{equation}\label{eq:theta-polynomial}
    \Theta_{a,b,c}(x)
    =x^{a+b+c}-x^a-x^b-x^c-2.
  \end{equation}

  \item If \(C\) is a dumbbell with cycle lengths \(a,b\geq1\) and corridor
  length \(c\geq1\), then
  \begin{equation}\label{eq:dumbbell-pressure}
    \alpha^{2c}(\alpha^a-1)(\alpha^b-1)=4,
  \end{equation}
  and \(\alpha\) is the unique root \(x>1\) of
  \begin{align}
    D_{a,b,c}(x)
    &=
    x^{a+b+2c}-x^{a+2c}-x^{b+2c}+x^{2c}-4.
    \label{eq:dumbbell-polynomial}
  \end{align}
\end{enumerate}
Conversely, the unique root \(x>1\) in each of these three families is the
Hashimoto spectral radius of the corresponding rank-two core.
\end{theorem}

\begin{proof}
The kernel classification preceding the theorem proves that the three cases
are exhaustive and mutually exclusive at the level of kernel type.

\emph{Figure-eight.}  For a figure-eight, the weighted turn
matrix commutes with the involution that exchanges the two orientations of
each loop.  Perron uniqueness therefore makes the positive Perron vector
invariant under this involution, so it takes the same
value on the two orientations of each loop.  Applying
\cref{thm:kernel-pressure} and writing the two reversal-invariant coordinates
gives \eqref{eq:figure-eight-pressure} directly.
Clearing denominators yields both \eqref{eq:figure-eight-product} and
\eqref{eq:figure-eight-polynomial}.  The left side of
\eqref{eq:figure-eight-pressure} is strictly decreasing from \(2\) to \(0\)
on \(x\in(1,\infty)\), so the root is unique.

\emph{Theta graph.}  For a theta graph, after traversing one of the three
paths, a non-backtracking path may return along either of the other two paths.
Put \(z=x^{-1}\).  The involution that interchanges the two branch vertices
and reverses all three oriented kernel edges commutes with the weighted turn
operator.  After normalization, Perron uniqueness makes its positive Perron
vector invariant under this involution.  The two orientations of each kernel
edge therefore carry the same Perron coordinate.  On the three reversal
pairs, the weighted turn equation is represented by the matrix with
off-diagonal entry \(z^{\ell_j}\), where
\((\ell_1,\ell_2,\ell_3)=(a,b,c)\), in column \(j\) and zero
diagonal.  The matrix determinant lemma, applied to
\[
  I+\operatorname{diag}(z^a,z^b,z^c)
  -\mathbf1_3(z^a,z^b,z^c),
\]
where \(\mathbf1_3\) is the three-dimensional all-ones column vector, shows
that its Perron root is one precisely when
\eqref{eq:theta-pressure} holds.  Indeed, under this scalar condition
the vector with coordinates \((1+z^{\ell_j})^{-1}\) is a strictly positive
eigenvector of eigenvalue one; irreducibility and Perron--Frobenius theory
then show that one is the spectral radius.  Clearing denominators gives
\eqref{eq:theta-polynomial}.  The left side of
\eqref{eq:theta-pressure} decreases strictly from \(3/2\) to zero.

\emph{Dumbbell graph.}  For the dumbbell, write
\(p=x^{-a}\), \(s_0=x^{-b}\), and \(w=x^{-c}\).
The weighted turn operator is invariant under independently exchanging
the two orientations of either loop.  Let \(A\) and \(B\) be the
Perron-vector values on the two loop-orientation pairs, and let \(U,V\) be the
values on the two corridor orientations.  The eigenvalue-one equations for
the weighted turn matrix are
\[
  (1-p)A=wU,\qquad V=2pA,\qquad
  (1-s_0)B=wV,\qquad U=2s_0B.
\]
Eliminating \(A,B,U,V\) gives
\[
  (1-p)(1-s_0)=4w^2ps_0,
\]
which is \eqref{eq:dumbbell-pressure}.  Its left side in the product form
\[
  x^{2c}(x^a-1)(x^b-1)=4
\]
is strictly increasing from zero to infinity for \(x>1\).  Expanding gives
\eqref{eq:dumbbell-polynomial}.  The converse in every case follows by
constructing the indicated subdivision and applying
\cref{thm:kernel-pressure}.
\end{proof}

\end{document}
